\chapter{数据保存协议与发布验证}

\section{三层数据结构}
\begin{enumerate}
  \item \textbf{Authoritative raw}: \code{runs_eval/revised/publication_seed101_v1/<network>/<controller>/<method>/}，保存不可变 manifest、environment、summary、result、压缩逐秒记录和 NPZ。
  \item \textbf{Validated release}: \code{output_result/revised/publication_seed101_v1/}，保存 release manifest、validation report、CSV、图表、paired bootstrap 和兼容 dashboard JSON。
  \item \textbf{Sanitized site bundle}: \code{docs/site/dist/data/publication/}，只保存字段白名单允许的聚合与脱敏 JSON。
\end{enumerate}

JSON 是权威格式，CSV 仅供阅读。JSON 禁止 NaN/Infinity；不适用或缺失值使用 \code{null}，并在 metric registry 说明原因、单位、方向和分母。Raw 数据不得复制到托管站点。

\section{Selection registry 与主键}
\code{runs_eval/revised/selections/final_evaluation_seed101.json} 冻结 8 个 parent、8 个 offline WCE、40 个 final continuation 和 40 个 planned evaluation suite。每项保存 repo-relative path、绝对本机路径、checkpoint step、目录聚合 SHA-256、manifest hash、manifest 文件 SHA-256、runtime/source compatibility 状态和计划输出目录。模型选择只允许使用该 allowlist，不允许按 mtime、“latest”或宽目录扫描推断。

Rollout 唯一键为
\begin{Verbatim}[fontsize=\small,breaklines=true]
network/controller/method/training_seed/split/scenario/
generation_seed/arrival_seed/sumo_seed/policy_seed/attempt
\end{Verbatim}
公开 rollout ID 使用上述稳定字段的摘要；绝对路径不进入公开 ID。

\section{Publication-complete 硬门槛}
只有全部规则同时通过，\code{publication_complete} 才能为 true：
\begin{itemize}
  \item accepted 恰好 9,200，rejected、missing、duplicate、extra 均为 0；
  \item 40 个 checkpoint 全部来自 selection allowlist，controller learning steps 为 2,320,000；
  \item 所有记录 \code{pilot=false}、\code{status=complete}、\code{horizon=sample_count=3600}；
  \item arrival 51001--51010 与 SUMO 61001--61010 按 index 精确对应；
  \item 同一 network/scenario/index 的 demand hash 与 SUMO seed 跨 20 个 controller-method 组合一致；
  \item protocol、scenario、profiles、network、artifact、checkpoint 和 manifest hashes 与冻结输入一致；
  \item 时间戳严格为 1--3600，queue 非负且 finite；
  \item $\mathrm{integrated\ queue}\approx3600\times\mathrm{mean\ queue}$，peak 不小于 mean；
  \item scheduled = inserted + pending，completed $\le$ inserted，completed-trip denominator = completed；
  \item NPZ 每车道 queue 之和与 JSONL network total 对齐；
  \item public bundle 不含绝对路径、token、checkpoint payload、raw log 或 pilot/history 数据。
\end{itemize}

\section{磁盘与压缩门槛}
2026-09-29 实测根文件系统约有 69 GiB 可用，已超过 60 GiB 启动阈值，但每次正式启动前都必须重新检查。按现有 pilot 估算，9,200 条 raw rollout 可能超过 33 GB。启动条件是“可用空间至少 60 GiB，且预计完成后保留至少 20 GiB”。若不满足，应先实施 suite-level shared provenance、精简 per-rollout manifest、gzip JSONL 和压缩 NPZ，并用 pilot 重新测量；不得通过删除未知用户数据解决。

\section{Reward 与环境验证}
发布验证还要计算 selected training runs 的 controller/WCE clipping fraction，生成 source/config/runtime compatibility matrix，并将 baseline 不同运行环境标为已知限制。环境差异不阻断用户选择的 release，但必须令 \code{strict_runtime_comparability=false}，并禁止训练速度或严格公平性的强结论。
